\documentclass[11pt]{article}
\usepackage[a4paper,margin=24mm]{geometry}
\usepackage[T1]{fontenc}
\usepackage{lmodern,microtype,amsmath,amssymb,amsthm,mathtools,booktabs,array}
\usepackage{xcolor}
\usepackage[colorlinks=true,linkcolor=blue!45!black,urlcolor=blue!55!black,citecolor=blue!45!black]{hyperref}
\usepackage{enumitem}
\hypersetup{pdftitle={Acquiring exact field counters from raw Boolean cubics},pdfauthor={Astra P3},pdfsubject={Constructive representations, rank-metric decoding, and exact counting}}
\setlist{nosep,leftmargin=*}
\setlength{\parskip}{4pt}
\setlength{\parindent}{0pt}
\newtheorem{theorem}{Theorem}[section]
\newtheorem{lemma}[theorem]{Lemma}
\newtheorem{corollary}[theorem]{Corollary}
\newtheorem{proposition}[theorem]{Proposition}
\theoremstyle{remark}\newtheorem{remark}[theorem]{Remark}
\newcommand{\F}{\mathbb F}
\newcommand{\Tr}{\operatorname{Tr}}
\newcommand{\rank}{\operatorname{rank}}
\newcommand{\poly}{\operatorname{poly}}
\newcommand{\im}{\operatorname{im}}
\newcommand{\Cent}{\operatorname{Cent}}
\newcommand{\wt}{\operatorname{wt}}
\newcommand{\supp}{\operatorname{supp}}
\newcommand{\softO}{\widetilde O}
\title{Acquiring exact field counters from raw Boolean cubics\\\large Hidden scalar fields, rank-metric decoding, and exact reconstruction}
\author{Astra P3 research construction}
\date{27 October 2026}
\begin{document}
\maketitle
\begin{abstract}
We construct an exact representation compiler for explicit Boolean cubics whose highest difference tensor reveals an extension field. The lower quadratic part may contain an unknown, dense error of small binary polar rank. The acquired field turns the compatible quadratic blocks into explicit Gabidulin codes. Decoding manufactures a field-polynomial base and the actual quadratic defect; a real Walsh expansion then reconstructs the original counts from exactly computable field counts. No field, coordinate transformation, code, defect or polynomial base is supplied.

For $n=bd$ original bits, Frobenius bandwidth $s\ge1$, and defect rank $r$ satisfying $2r\le d-2s-1$, the resulting unconditional bit-time bound is $2^{r+O(b(\log n+s))}\poly(L+n)$. In particular, fixed $b$, ordinary field degree three and $r=O(\log n)$ give polynomial time and space. The general evaluator is the established Lauder--Wan point-counting algorithm. A complete narrower counter, including an actual finite Dwork trace computation, is implemented and tested. An exact symbolic front end extends it to accepted gate programs while charging all intermediate work against their original input count. This document distinguishes new interface claims from existing centroid, decoding, Gaussian and point-counting machinery; it makes no first-priority, generic-circuit, or practical-speed claim.
\end{abstract}

\section{Problem, scope, and the obstruction removed}
The input is a Boolean algebraic normal form
\[
 p(x)=\bigoplus_{S\subseteq[n],\ |S|\le3} a_S\prod_{i\in S}x_i,
 \qquad x\in\F_2^n.
\]
Let $m$ be its number of nonzero monomials. A list of their at most three indices has length $L=O(n+(m+1)\log(n+2))$ bits; the programs instead use integer bit masks, costing at most $O(n(m+1))$ bits. All bounds are polynomial in either honest encoding length plus $n$. The independent input count remains $n$, even when $m=\Theta(n^3)$.

We seek the signed count and both output counts:
\[
 Z(p)=\sum_{x\in\F_2^n}(-1)^{p(x)},\qquad
 N_0=\frac{2^n+Z(p)}2,\quad N_1=\frac{2^n-Z(p)}2.
\]
An invertible acquired coordinate map preserves multiplicities exactly. No Tseitin variables are charged as original inputs, and no conversion to a linear-size circuit or CNF is used.

The structural family is restricted. Its third difference tensor must reveal a useful field, and its lower quadratic part must admit the stated bounded-distance decoding. These are conditions under which the compiler is guaranteed to find a certificate, not a supplied decomposition. On an arbitrary input, every successful output is an exact identity; failed structural tests reject the route. Nothing here is a polynomial-time solver for arbitrary circuits or SAT.

This is an alternative to extending the governing replacement-form manuscript directly through unrestricted program states. The pivot is justified by a concrete obstruction: a dense lower quadratic can conceal a tractable field computation, even after the top tensor reveals its field. Fitting the entire quadratic as field-compatible fails as soon as a noncompatible error is present; asking for the error as input would assume the useful representation. The decoding step removes that assumption within a quantified, adversarial error radius.

The main theorem below is independent of the pinned ETH manuscript \cite{eth}. Two other acquired representation constructions that still use the earlier local-phase theorem are preserved in the source archive and stated separately in Section~\ref{sec:other}. Their dependency is not transferred to the unconditional field theorem.

\section{The acquired-field counting theorem}
For a cubic $p$, put $T(u,v,w)=\Delta_u\Delta_v\Delta_w p(0)$, where $\Delta_u p(x)=p(x+u)+p(x)$. This is an alternating symmetric trilinear form over $\F_2$. Its radical consists of $u$ such that $T(u,v,w)=0$ for every $v,w$. Define
\[
 \Cent(T)=\{A\in\operatorname{End}_{\F_2}(V):
 T(Au,v,w)=T(u,Av,w)=T(u,v,Aw)\}.
\]
Assume $T$ has zero radical and that this computed algebra is a field $F$ of binary dimension $d$. Then $V$ is an $F$-module of dimension $b=n/d$. The field presentation and module coordinates are acquired, as proved below.

In any such coordinates, the trace pairing constructs an alternating $F$-trilinear form $\mathcal T$ with $T=\Tr_{F/\F_2}\mathcal T$. Set
\[
 C(X)=\sum_{i<j<k}\mathcal T(e_i,e_j,e_k)X_iX_jX_k.
\]
The function $p-\Tr C$ is quadratic. Write its alternating polar matrix in $b\times b$ blocks of $d\times d$ bits. Each block has a unique expression
\[
 B_{ij}(u,v)=\Tr\bigl(uL_{ij}(v)\bigr).
\]
For an integer $s\ge1$ with $2s<d$, let $\mathcal B_s$ be the space of symmetric alternating block matrices whose maps have the form
\begin{equation}\label{eq:band}
 L_{ij}(X)=\sum_{a=-s}^{s}c_{ij,a}X^{2^a}.
\end{equation}
Negative exponents here mean inverse Frobenius, $X^{2^{-a}}=X^{2^{d-a}}$ as functions on $F$. They do not mean negative ordinary polynomial powers.

\begin{theorem}[Acquired rank-corrected field count]\label{thm:main}
Suppose the raw input has the centroid and radical properties above. Suppose its quadratic polar matrix satisfies
\[
 B=B_0+D,\qquad B_0\in\mathcal B_s,\qquad
 \rank_{\F_2}D=r,\qquad 2r\le d-2s-1.
\]
Neither $B_0$ nor $D$ is supplied. A deterministic polynomial-time acquisition algorithm recovers both matrices, an explicit coordinate map, and a field polynomial $P_0$ of ordinary degree at most
\[
 D_s=2^s+1
\]
such that $p(x)=\Tr P_0(X(x))+q_D(X(x))$, where $q_D$ is a canonical binary quadratic with polar matrix $D$. It then computes $Z(p),N_0,N_1$ in bit time
\begin{equation}\label{eq:explicitcost}
 \poly(L+n)+2^r\softO\!\left(
 2^{O(v)}d^{3v+7}v^{3v+5}D_s^{3v}2^{2v+4}\right),
 \qquad v=b+1.
\end{equation}
A simpler uniform upper bound is
\[
 2^{r+O(b(\log(n+2)+s))}\poly(L+n).
\]
The corresponding space bound is
\[
 \poly(L+n)+\softO\!\left(
 2^{O(v)}d^{2v+4}v^{2v+3}D_s^{2v}2\right).
\]
Walsh leaves are streamed, so no $2^r$ space factor is required.
\end{theorem}

\begin{corollary}[Ordinary field cubics]
Suppose, over the acquired centroid field, some hidden affine field coordinates satisfy
\[
 p=\Tr P+e,\qquad \deg_F P\le3,
 \qquad \rank(\operatorname{polar}(e))=r,
 \qquad 2r\le d-3.
\]
Then Theorem~\ref{thm:main} applies with $s=1$. For fixed $b$ the time is $2^r\poly(L+n)$ and space is polynomial. In particular, $r=O(\log n)$ gives polynomial exact counting; $r=o(n)$ gives subexponential exact counting. The quadratic error may be dense and is acquired from the raw input.
\end{corollary}

The model must be over the acquired centroid field. Low degree over an unspecified proper subfield need not remain low degree over a larger field. An affine change over the correct field is harmless: an ordinary field cubic stays cubic, and its repeated-variable cubic terms contribute precisely the bandwidth-one quadratic blocks.

For growing parameters, the subexponential condition is
\[
 r+b(\log(n+2)+s)=o(n),
\]
together with the decoding radius. Fixed $b$ and $r,s=O(\log n)$ give polynomial time and space. Fixed $b$ and $r,s=o(n)$ give subexponential time, but the published field evaluator need not use polynomial space when $s$ grows faster than $\log n$.

No useful rank or bandwidth need be guessed correctly by a user. After the field is acquired, one may try every $1\le s\le\lfloor(d-1)/2\rfloor$, retain only verified decompositions, and select a candidate minimizing the explicit leaf-cost bound. There are only $O(d)$ trials, each polynomial-time. For fixed acquired $b,d$, comparing $2^{r_s}D_s^{3v}$ uses integers of polynomial bit length. This selection rule is also implemented and tested: a width-one trial failed on the tested wider fixture, width two recovered rank zero, and width three was valid but more expensive. The delivered specialized count backend fixes $s=1$; general acquisition accepts either an explicit bandwidth or this automatic selection.

\section{Manufacturing the field and the canonical cubic}
\subsection{Centroid equations and field certification}
The entries of $T$ are read directly from the cubic coefficients. Its radical is the kernel of a binary matrix. The centroid equations are linear in the $n^2$ entries of $A$, with $O(n^3)$ equations. Binary elimination computes a basis in polynomial time. The implementation sometimes stops adding tensor slices early, but verifies every remaining equation against the candidate kernel basis before accepting it.

For completeness, the centroid is commutative when the radical is zero. Moving $A,B$ between different tensor slots gives both
\[
 T(ABu,v,w)=T(u,BAv,w),\qquad
 T(ABu,v,w)=T(u,ABv,w).
\]
Nondegeneracy therefore gives $AB=BA$. The same transfers prove closure under multiplication, and the identity belongs to the algebra.

On a finite-dimensional commutative binary algebra, the square map $\sigma:a\mapsto a^2$ is linear. Injectivity of $\sigma$ is equivalent to reducedness: a nonzero nilpotent would produce a nonzero element killed by squaring after sufficiently many squarings. A reduced finite algebra is a product of finite fields. Its $\sigma$-fixed space has one binary dimension for each field factor. Thus
\[
 \rank\sigma=d,\qquad \dim\ker(\sigma-I)=1
\]
certify that the centroid is one field. These tests are polynomial-size linear algebra, not a search over its $2^d$ elements.

\subsection{A deterministic field generator without enumerating the field}
Scanning a basis for an element of degree $d$ often succeeds, but it is not the justification for the complexity bound. A deterministic fallback is as follows. Factor the integer $d$ into pairwise coprime prime powers $m_i=\ell_i^{e_i}$, using trial division polynomial in $d$. In the field algebra compute
\[
 K_i=\ker(\sigma^{m_i}-I),\qquad
 K'_i=\ker(\sigma^{m_i/\ell_i}-I).
\]
Choose a basis vector $\gamma_i\in K_i\setminus K'_i$. Its field degree is exactly $m_i$. The sum $\alpha=\sum_i\gamma_i$ has degree $d$.

To verify the last claim, let $\gamma,\eta$ have coprime degrees $m,t$. If $\sigma^j$ fixes their sum, then
\[
 \sigma^j\gamma-\gamma=\sigma^j\eta-\eta=\delta
 \in \F_{2^m}\cap\F_{2^t}=\F_2.
\]
At least one of $m,t$ is odd. Iterating the corresponding Frobenius identity that many times gives $\delta=0$. Hence $\sigma^j$ fixes both elements, and the sum has degree $mt$. Induction proves the claim. For $d=1$ take the identity element.

The first relation among $1,\alpha,\ldots,\alpha^d$ supplies an irreducible polynomial $h$ of degree $d$ and the presentation $F=\F_2[t]/(h)$. A greedy module-basis construction adds an original binary basis vector outside the span already obtained, together with its $d$ field multiples. Field multiplication makes each added orbit independent. This yields an invertible $n\times n$ binary coordinate matrix. Every step is polynomial in the dimensions of explicit matrices. Here the centroid is the balanced input-slot algebra explicitly defined above; the scalar trace lift is proved separately. Hidden-field descent has established tensor ancestry, including the documented \texttt{TensorOverCentroid} operation \cite{magma}; the argument here exposes the costs used by this compiler.

\subsection{Trace lift and exactness}
The trace pairing $(a,z)\mapsto\Tr(az)$ is nondegenerate. Define $\mathcal T$ by
\[
 \Tr(a\mathcal T(u,v,w))=T(au,v,w)\quad(a\in F).
\]
Centroid identities make $\mathcal T$ field-linear in all slots. Alternation is also inherited, since
\[
 T(au,u,w)=T(u,u,aw)=0
\]
for all $a$, so nondegeneracy of the pairing forces $\mathcal T(u,u,w)=0$. Its coefficients are obtained by solving $d\times d$ binary trace Gram systems. The squarefree field cubic $C$ defined earlier therefore has precisely the original third tensor after trace. The residual is a Boolean quadratic as an exact coefficient identity.

An affine field-coordinate change can add repeated-variable terms such as $X_i^2X_j$ to this canonical lift. Such terms are binary quadratic. They must be included in the lower representation: treating the canonical lift as the unique ordinary field cubic would lose valid computations.

\section{Manufacturing and decoding the quadratic code}
\subsection{The code is an output of acquisition}
Let $g_1,\ldots,g_d$ be the binary power basis of the acquired field and $G$ its trace Gram matrix. Each binary matrix block is converted to a linear map $L:F\to F$ by applying $G^{-1}$ to its columns. This is an invertible change of output coordinates and preserves binary rank.

For a compatible block in \eqref{eq:band}, precompose by the invertible Frobenius map:
\[
 f(X)=L(X^{2^s})=\sum_{j=0}^{2s}c_{j-s}X^{2^j}.
\]
The evaluations $(f(g_1),\ldots,f(g_d))$ form a length-$d$, dimension-$k$ Gabidulin code over $F$, where $k=2s+1$. A nonzero linearized polynomial of Frobenius degree at most $k-1$ has at most $2^{k-1}$ roots. Its kernel dimension is at most $k-1$, so its binary rank is at least $d-k+1=d-2s$. This proves the required minimum-distance bound directly; the code and its decoders have established ancestry \cite{silva,bombar}.

Every block of the actual defect $D$ has rank at most $r$, since block extraction is a row and column restriction. Trace-coordinate conversion and Frobenius precomposition preserve that rank. Thus each received block is uniquely decodable at radius
\[
 t=\left\lfloor\frac{d-k}{2}\right\rfloor
 =\left\lfloor\frac{d-2s-1}{2}\right\rfloor.
\]
The field, evaluation points and code coefficients were all manufactured from the raw tensor. There is no assumed free code decomposition.

\subsection{An explicit deterministic interpolation decoder}
Suppose the received map is $y=f+e$, with $\deg_2 f<k$ and $\rank e\le t$. Here $\deg_2$ is the largest index in a linearized polynomial, not its ordinary degree. Solve for a monic $E$ and $N$ satisfying
\begin{equation}\label{eq:decoder}
 E(y(g_j))=N(g_j)\quad(1\le j\le d),\qquad
 \deg_2 E=t,\quad \deg_2N<k+t.
\end{equation}
Writing $E(X)=X^{2^t}+\sum_{i<t}e_iX^{2^i}$ makes this a linear system over $F$ in $k+2t\le d$ unknown coefficients.

There is a solution. The image of $e$ is a binary subspace of dimension $u\le t$ and has a monic linearized annihilator of Frobenius degree $u$. Composing with enough squarings gives a monic annihilator of degree $t$. Choose $N=E\circ f$.

Every solution of \eqref{eq:decoder} works. The polynomial $H=N-E\circ f$ vanishes on $\ker e$, whose dimension is at least $d-t$. But $\deg_2 H\le k+t-1<d-t$. The root bound forces $H=0$. Recover $f$ by composition division with $E$ on the left: from a leading coefficient $c$ at index $t+j$, recover the next coefficient $a$ by $a^{2^t}=c$, subtract $E\circ(aX^{2^j})$, and continue. Frobenius is invertible, so this step is exact and polynomial-time. Finally verify zero remainder and rank distance at most $t$.

This is the established left-annihilator decoding formulation; confusing it with a right-hand interpolation system would incorrectly assert extension-field linearity. The implementation and an independent audit both check the orientation. Errors may be adversarial and may have rank strictly less than the chosen maximal radius.

There are $b^2$ blocks. Their interpolation systems have at most $d$ rows and unknowns over a degree-$d$ binary field. Ordinary Gaussian elimination uses $O(b^2d^3)$ field operations. With classical binary polynomial arithmetic this is polynomial bit cost, for example $O(b^2d^5)$ for the dominant multiplications. Undo Frobenius, reassemble $B_0$, verify symmetry and zero diagonal, and compute the exact rank of $D=B+B_0$. Under the promise the planted matrices are recovered. Outside it, only a verified identity is accepted.

\section{Integrating the decoded blocks and reconstructing counts}
\subsection{A polynomial of bounded ordinary degree is constructed}
For $i<j$, a positive-index block term gives
\[
 \Tr(c_aX_iX_j^{2^a}),\qquad 0\le a\le s.
\]
A negative-index term is folded without increasing the promised degree beyond $2^s+1$:
\[
 \Tr(c_{-a}X_iX_j^{2^{-a}})
 =\Tr(c_{-a}^{2^a}X_i^{2^a}X_j),\qquad 1\le a\le s.
\]
For a diagonal block, self-adjointness of $L$ imposes
\[
 c_{-a}=c_a^{2^{d-a}}.
\]
These paired terms contribute zero to $\Tr(XL(X))$. Alternation therefore forces $c_0=0$. The quadratic refinement
\[
 Q_i(X_i)=\Tr\left(\sum_{a=1}^{s}c_aX_i^{1+2^a}\right)
\]
has exactly the required diagonal polar form. The strict inequality $2s<d$ keeps every positive and negative index distinct.

The refinement need not equal the canonical binary quadratic: its values on individual binary basis vectors may be nonzero. Correct this with a Boolean linear function, represented uniquely by a field-trace linear form through $G^{-1}$. More generally, after adding the canonical cubic $C$, match the raw values at zero and at all binary basis vectors. This supplies all remaining affine terms. For a constant one, choose an acquired field element of trace one; taking the field element $1$ is wrong when $d$ is even.

The resulting field polynomial $P_0$ has ordinary degree at most $D_s=2^s+1$. It uses $O(b^3+b^2s)$ sparse monomials. With its coordinate matrix and a rank factorization of the defect, the representation occupies
\[
 O\bigl(n^2+nr+d(b^3+b^2s)\bigr)
\]
bits up to ordinary index-encoding factors. This is manufactured from the raw coefficient list. Equality follows from equality of the third tensor, quadratic polar form and affine part. Alternatively, every constructed trace monomial has Boolean degree at most three, so checking all coefficients by evaluations at Hamming weights at most three is a polynomial-size global identity test.

The family is invariant under affine changes of its $b$ field coordinates for $s\ge1$. A Frobenius power of an affine field-linear form is again affine in the corresponding Frobenius powers. A quadratic term with index difference in $[-s,s]$ stays in that range after substitution and trace folding. The transformed ordinary cubic contributes only additional bandwidth-one terms. A different field presentation conjugates coefficients and preserves the trace.

\subsection{Real Walsh reconstruction, with exact arithmetic}
The rank $r$ of an alternating binary matrix is even. Exact quadratic elimination gives independent binary linear forms $\ell_1,\ldots,\ell_r$ and an affine correction such that
\[
 q_D(y)=\sum_{j=1}^{r/2}\ell_{2j-1}(y)\ell_{2j}(y)+\ell_0(y).
\]
The canonical $q_D$ has zero constant, so $\ell_0$ can be chosen linear. For two bits $u,v$,
\[
 (-1)^{uv}=\frac12\sum_{\alpha,\beta\in\F_2}
 (-1)^{\alpha\beta+\alpha u+\beta v}.
\]
Multiplying these identities yields exactly $2^r$ real-character leaves:
\begin{equation}\label{eq:walsh}
 Z(p)=2^{-r/2}\sum_{a\in\F_2^r}\epsilon_a
 \sum_{y\in\F_2^n}(-1)^{\Tr P_0(X(y))+\ell_0(y)+\sum_i a_i\ell_i(y)}.
\end{equation}
Every binary parity in a leaf is a field-trace linear form. Thus each leaf has the same field-variable count and ordinary degree bound. No complex or quarter phases are introduced.

Stream selectors $a$, evaluate one leaf, and accumulate a signed integer numerator. Its absolute value is at most $2^{n+r}$, so $O(n+r)$ bits suffice. Divide exactly by $2^{r/2}$ after all leaves, then obtain $N_0,N_1$ by the displayed integer formulas. The public complete counter streams by default. Retaining every leaf certificate is an explicitly optional finite-experiment mode with exponential storage, not the polynomial-space algorithm.

\subsection{Exact field evaluation and its charged cost}
For a leaf polynomial $G$ on $F^b$, construct
\[
 W^2+W=G(X_1,\ldots,X_b).
\]
For each $X$, this Artin--Schreier equation has two field solutions precisely when $\Tr G(X)=0$. If its affine point count is $N_{\rm AS}$, then
\[
 N_{\rm AS}=2N_0(G),\qquad Z(\Tr G)=N_{\rm AS}-2^n.
\]
The hypersurface has $v=b+1$ variables and ordinary degree at most $D_s$. Lauder--Wan Proposition~36 and the torus decomposition in their Section~6.4 give the bit-time and space terms in Theorem~\ref{thm:main}, with characteristic $2$, field extension degree $d$, and extension-count index $1$ \cite{lauderwan}. No nonsingularity assumption is needed.

Only one point count is required. Their full-zeta Theorem~37 computes many extension counts and incurs an unnecessary larger dependence on $b$. For this task the modulus precision is $d(v+1)=O(n+d)$ bits. The cited costs include unramified-ring arithmetic; the trace-to-count conversion is just integer subtraction. Processing the torus pieces sequentially avoids multiplying space by their number.

The field degree bound is not a free compression of computational work. Since $\log_2D_s=\Theta(s)$, its contribution to the exponent is $O(bs)$. A succinct high-degree polynomial or program cannot be charged as a low-degree one merely because its formula is short. These costs complete the proof of Theorem~\ref{thm:main}.

\section{A complete executed counter for hidden triple products}\label{sec:complete}
The delivered complete counter handles $b=3$, $s=1$, and the subfamily whose recovered ordinary homogeneous field cubic is the product of three independent field-linear forms. This includes every binary affine disguise of
\[
 \Tr\bigl(\lambda XYZ+A XY+B XZ+C YZ+D X+E Y+F Z+G\bigr)+e,
 \quad \lambda\ne0,
\]
with an unknown quadratic defect within the decoding radius. It does not execute the general multivariate Lauder--Wan backend. Its narrower leaf evaluator is fully implemented, including the finite trace computation described below.

\subsection{Acquiring the ordinary cubic and its factor planes}
The top tensor alone gives the canonical $\Tr(\lambda XYZ)$. Under an unknown $\operatorname{GL}_3(F)$ change, however, a product cubic also acquires $X_i^3$ and $X_i^2X_j$ terms. These are only binary quadratic and are invisible to the third difference. The implemented fitter therefore includes all three cross quadratics, three pure cubes, six ordered $X_i^2X_j$ terms, field-linear terms and a trace-one constant. It solves an exact binary coefficient system after replacing the raw quadratic matrix by the decoded matrix. The generic bandwidth implementation instead integrates the decoded blocks directly.

For $d\ge3$, the homogeneous ordinary degree-three part is uniquely determined by the scalar trace function modulo lower ordinary field degree. Indeed, expand the trace as a sum of Frobenius conjugates and reduce positive exponents modulo $2^d-1$. The exponent patterns $(3)$, $(2,1)$ and $(1,1,1)$ each have a length-$d$ Frobenius orbit whose only member of total ordinary degree at most three is the original pattern. No lower-degree pattern belongs to those orbits. Equality of scalar trace functions therefore fixes every homogeneous cubic coefficient. Linear/square trace aliases and trace-zero constants remain possible and are handled explicitly.

Suppose this homogeneous cubic is $H=cL_1L_2L_3$ with independent $L_i$. Test seven distinct field elements in $v(t)=(1,t,t^2)$. Each $L_i(v(t))$ is a nonzero polynomial of degree at most two, so $H(v(t))\ne0$ for one probe. With this $v$, test seven vectors $u(s)=(1,s,s^2)$. The three roots of $H(u+Tv)$ are distinct when
\[
 \prod_{i<j}\bigl[L_j(v)L_i(u)+L_i(v)L_j(u)\bigr]\ne0.
\]
This is a nonzero polynomial of degree at most six in $s$. A squarefree slice is therefore found in seven probes, using a gcd with the derivative.

Verify that the slice divides $T^{2^d}-T$ by $d$ modular squarings. Split it deterministically using gcds with $\Tr(aT)$ modulo the slice, for a binary basis of field elements $a$. The trace pairing separates its three roots. At a root $r_i$, the gradient $\nabla H(u+r_iv)$ is a nonzero scalar multiple of the corresponding factor normal. Invert these three normals and verify the entire transformed homogeneous cubic is exactly a nonzero multiple of $XYZ$. These operations are polynomial in $d$; there is no $2^d$-element search. Inputs with a nonfactorable cubic are rejected by this specialized route, even though the general theorem may apply.

\subsection{Reduction to a Kloosterman sum}
After this field-linear change, fold field-square terms using
\[
 \Tr(cX^2)=\Tr(c^{2^{d-1}}X).
\]
This step is essential: omitting it caused a failed raw-input test during construction and was repaired before the reported executions.

Scale $X$ to make its cubic coefficient one. With normalized coefficients $A,B,C,D,E,F,G$, translate $X,Y,Z$ to remove the three cross quadratics. The result is
\[
 XYZ+aX+bY+cZ+g,
\]
where
\[
 a=D+AB,\quad b=E+AC,\quad c=F+BC,\quad
 g=G+DC+EB+FA.
\]
Let $q=2^d$ and $\psi(z)=(-1)^{\Tr z}$. Summing over $Z$ gives
\[
 q\psi(g)\sum_{XY=c}\psi(aX+bY).
\]
For $c=0$, the inner sum is $q\mathbf1_{a=0}+q\mathbf1_{b=0}-1$. For $c\ne0$, it is $q-1$ if $a=b=0$, and $-1$ if exactly one of $a,b$ vanishes. Otherwise it is
\[
 K(abc),\qquad K(t)=\sum_{x\in F^*}\psi(x+t/x).
\]
Thus only an exact one-variable Kloosterman evaluator remains. Walsh characters change only the three field-linear coefficients, so every reconstruction leaf stays in this family.

\subsection{An actual finite Dwork trace, rather than field enumeration}
The implementation specializes established Dwork machinery \cite{lauderwan}. Let
\[
 R_N=(\mathbb Z/2^N\mathbb Z)[t]/(\widetilde h(t)),
\]
where $\widetilde h$ lifts the acquired irreducible binary modulus. This is an unramified ring presentation. Compute the Teichmuller lift $A$ of a nonzero Kloosterman parameter by repeated $2^d$ powers. After $d\lceil N/d\rceil$ squarings the lift is correct modulo $2^N$; the code checks its residue and $A^{2^d}=A$.

Use the formal splitting series
\[
 \theta(z)=\exp(-2z+2z^2)=\sum_{j\ge0}\lambda_jz^j,
 \qquad
 j\lambda_j=-2\lambda_{j-1}+4\lambda_{j-2}.
\]
Its coefficients satisfy $\operatorname{ord}_2(\lambda_j)\ge j/4$. Compute them as exact rational numbers, whose reduced denominators are odd, and reduce modulo $2^N$. Terms with $j\ge4N$ vanish. A subtlety is important: $\theta(1)$ is the convergent coefficient-series value $-1$; substituting into the formal expression and writing $\exp(0)=1$ is invalid at that boundary. The code verifies the coefficient sum modulo $2^N$.

For $i=0,\ldots,d-1$, let
\[
 F_i(X)=\theta(X)\theta(A^{2^i}/X),\qquad
 M_i[u,v]=[X^{2u-v}]F_i(X),
\]
with $u,v$ ranging from $-(4N-1)$ to $4N-1$. These are matrices of dimension $8N-1$ over $R_N$. The trace identity is
\begin{equation}\label{eq:dwork}
 K(t)\equiv(q-1)\operatorname{Tr}(M_{d-1}\cdots M_0)\pmod{2^N}.
\end{equation}
One way to see it is to expand the splitting character over the Teichmuller roots of unity. Their sums annihilate every Laurent exponent not divisible by $q-1$. A cyclic matrix path has exponents $2u_{i+1}-u_i$, whose $2^i$-weighted sum is $(q-1)u_0$, exactly the surviving condition.

The finite truncation is exact modulo $2^N$. A Laurent coefficient of exponent $a$ in $F_i$ has valuation at least $|a|/4$. If a closed path contains an index of maximum absolute value at least $4N$, the edge entering that index has $|2u-v|\ge4N$ and vanishes. Hence omitted paths contribute zero. The Frobenius product in \eqref{eq:dwork} is a scalar integer residue; the implementation verifies that all nonconstant field-basis coefficients vanish.

The elementary bound $|K(t)|\le q-1$ makes $N=d+2$ sufficient for centered integer reconstruction. The executed triple counter uses the optional Hasse bound $|K(t)|\le2\sqrt q$, allowing $N=\lfloor d/2\rfloor+3$. For $t\ne0$, this follows from the nonsingular binary elliptic curve $y^2+xy=x^3+t$, whose point count is $q+1+K(t)$, using invariance of $K$ under Frobenius. Either precision choice is polynomial in $d$.

With classical multiplication a conservative bound for this specialized evaluator is $O(d^8)$ bit operations and $O(d^4)$ bits of space. The $d$ matrices have $O(d)$ rows; a ring element has $d$ coefficients of $O(d)$ bits. Products are formed sequentially. NumPy modular matrix products accelerate small precisions; when $N\le62$, reduction after 64-bit wraparound is valid because $2^N$ divides $2^{64}$. For larger $N$ the code uses arbitrary-precision objects. Public modulus inputs are checked for irreducibility. None of these implementation choices changes the exact arithmetic contract.

\section{Executed evidence and what it verifies}\label{sec:evidence}
The programs use Python 3 and NumPy for the finite trace products. Fixture construction knows a planted description only to generate raw coefficients and independent checks. The public compiler receives only the raw ANF and $n$. It does not read planted coordinates or reference counts.

\subsection{Actual counters}
The elementary Dwork backend matched a separate exhaustive field-sum oracle in 20 cases over degrees one through six. The full acquired field-only counter ran on raw hidden affine fixtures of 9, 12, 15 and 18 bits. The first three cases had full Boolean count and representation checks, totaling 37,376 pointwise identities. Larger fixtures at 24, 30 and 48 bits had 1,024 extra original-coordinate identity checks each and exact field-counter reference counts.

\begin{center}
\begin{tabular}{rrrr}
\toprule
Original bits & Acquired defect rank & Signed count $Z$ & Zero count $N_0$\\
\midrule
21 & 2 & $128$ & $1\,048\,640$\\
24 & 2 & $2\,304$ & $8\,389\,760$\\
30 & 2 & $35\,840$ & $536\,888\,832$\\
36 & 4 & $-176\,128$ & $34\,359\,650\,304$\\
\bottomrule
\end{tabular}
\end{center}
These four raw defect fixtures recovered the exact planted polar defect after coordinate conversion. The first two signed counts were independently verified over all $2^{21}$ and $2^{24}$ Boolean assignments. Each fixture had 1,024 further representation and quadratic-factor identities. The 30- and 36-bit cases used the actual exact Dwork leaves, not exhaustive Boolean leaves; their whole truth cubes were not enumerated. Forty-two separate root decoding cases tested all ranks through the radius in their selected fields.

A field-only 96-bit raw fixture contained 73,822 ANF terms, including 71,501 cubic terms. The compiler acquired $d=32$, $b=3$ and returned
\[
 Z=-302\,919\,748\,419\,584,
\]
\[
 N_0=39\,614\,081\,257\,132\,017\,336\,897\,765\,376.
\]
No $2^{96}$ Boolean cube or $2^{32}$ field loop was used. Its independent planted reference reduces to $K_{32}(1)=-70\,529$, checked through the elliptic-curve recurrence
\[
 K_j(1)=-K_{j-1}(1)-2K_{j-2}(1),\quad K_0(1)=-2,\ K_1(1)=1.
\]
There were 256 additional original/acquired-coordinate identity checks. This large example has defect rank zero; it must not be presented as a 96-bit corrupted-field experiment. Recorded local wall times are diagnostics, not benchmark or practical-speed claims.

\subsection{Wider representation and independent audits}
The general bandwidth compiler recovered five raw representations with $(n,b,s,r)$ equal to
\[
 (21,3,2,0),\ (27,3,2,2),\ (33,3,3,2),\
 (35,5,1,2),\ (45,5,2,2).
\]
Every planted defect matrix was recovered exactly. All 5,120 extra pointwise checks passed; the 21-bit case also had a complete cubic coefficient identity check. These executions validate the wider representation. They do not execute its general field point-counter.

An independent decoder and reconstruction implementation passed 144 Gabidulin cases, 144 block decodings across 16 fixtures, 1,024 real Walsh identities and 56 quadratic Gaussian sums checked exhaustively. Its base quadratic matrices had ranks up to 32, while errors of rank at most four had as many as 270 scalar edges. This tests an unknown low-rank error against a high-rank background.

A separate bandwidth audit checked all 512 maps over $F_8$ and all 4,096 maps over $F_{16}$ at bandwidth one; 1,081,344 cross-block folding identities; 69,632 precomposition identities; all 5,120 positive diagonal coefficient tuples for $(d,s)=(5,2),(6,2)$; 17,830,400 diagonal polar identities including the necessary linear correction; 29,456 half-period boundary checks; and 64 decoding cases through $s=4,d=13$. These are finite exact validations of the algebra, not a substitute for its proof.

The source archive contains all fixture generators, raw inputs, recovered maps, exact outputs and independent audit programs. Programs that retain leaf certificates for finite experiments identify that storage choice. The public complete counter was separately checked to process four leaves while retaining none, with the same exact count.

\section{A direct gate-program interface with charged intermediate states}\label{sec:program}
The construction can begin with an actual straight-line Boolean circuit rather than a supplied ANF. Let the circuit have $n$ original inputs, $S$ internal gates of types XOR, AND, NOT and constants, and one designated output. Wire indices give description length $|C|=O((n+S)\log(n+S+2))$ bits. The following front end is ordinary exact symbolic propagation; it is not a new circuit-equivalence algorithm.

At each wire, manufacture its complete Boolean ANF as a set of monomial masks. XOR is symmetric difference, NOT toggles the empty monomial, and AND takes all unions of two masks, canceling equal products modulo two. Multilinearity is imposed exactly through mask union. Impose an explicit term budget $M$ only \emph{after} all cancellation at a gate. Reject if the surviving ANF has more than $M$ terms. High-degree sparse intermediate states are allowed. At the final wire, verify that the actual output ANF has degree at most three and feed it to the field compiler.

\begin{proposition}[Certified program-to-field composition]
Every successful front-end execution proves its output ANF equal to the original circuit on every input. With a term budget $M$, sorting the at most $M^2$ products at each AND gives a conservative bit-time bound
\[
 O\bigl(|C|+n^2+SM^2n\log(M+2)\bigr)
\]
and storage $O(|C|+n^2+(SM+M^2)n)$ bits. The bit cost of the subsequent acquired field counter is added to this preprocessing cost. Internal wires contribute to $S$, not to the independent count $n$.
\end{proposition}
\begin{proof}
Induct on topological gate order. The Boolean monomial identity $x^Ax^B=x^{A\cup B}$ makes each computed gate polynomial exact, including all cancellations. Both operands have at most $M$ terms if execution has reached the gate. Sorting and parity cancellation therefore examine at most $M^2$ masks of at most $n$ bits. All stored wire polynomials and one product list fit the stated space bound. Explicit input masks contribute $O(n^2)$ bits. The output-degree check is performed on the exact polynomial, so it is a certificate rather than a sampling assumption.
\end{proof}

The implemented default is $M=\sum_{j=0}^3\binom nj=O(n^3)$. Thus any successful gate compilation has polynomial overhead in $n+S$. The budget is charged as a numerical parameter; an arbitrarily large binary-encoded budget is not free. This is still a restricted route: an intermediate ANF can exceed the budget even if later gates cancel it completely. The circuit step neither assumes a hidden field nor supplies its coordinates to the next stage.

Two complete gate-only fixtures were executed. The 21-input circuit has 879 internal gates, including 197 ANDs; the 36-input circuit has 3,205 gates, including 578 ANDs. Their maximum intermediate ANF sizes were 852 and 4,023 respectively. They implement hidden field computations and unknown rank-two or rank-four quadratic defects without giving those descriptions to the compiler. The resulting raw ANFs matched the separately preserved raw fixtures coefficient-for-coefficient, and the acquired exact counts were $Z=128$ and $Z=-176\,128$, respectively. Each had 1,024 extra original-input gate-evaluation checks. These are actual program executions with many internal states, not timing comparisons against SAT or other competitors.

A further 886-gate, 21-input fixture contains a degree-six intermediate monomial which is mixed into the output and then cancels. The term-budget front end accepted it and the complete counter again returned $Z=128$. This repairs the unnecessarily restrictive first implementation, which rejected any high-degree intermediate. Separately, a six-input program forming $\prod_i(1+x_i)$ and then XORing it with itself was rejected at the default 42-term budget, but accepted at budget 64 with exact zero output. The failure is an honest resource boundary of this symbolic representation.

The independent gate audit checked 580 quotient-ring identities and 160 generated circuits. Its term-budget compiler accepted 150 circuits, all checked against direct gate execution on 9,600 assignments. It also accepted sparse cancellation examples with intermediate degrees 4, 8 and 20 at term budget one. Full source and exact data accompany this proof.

\subsection{Why low-degree interpolation cannot replace certification}
Truncating ANFs above degree three does preserve every lower-degree coefficient: the span of higher-degree monomials is an ideal in the Boolean polynomial ring, since union multiplication never shrinks support. The unsound inference is to treat the truncated polynomial as the whole function. The four-input AND circuit has truncated cubic polynomial zero, yet its true signed count is $14$, not $16$. The executable front end rejects its noncubic final output.

There is also a precise general obstruction.
\begin{proposition}[Final-degree recognition is coNP-complete]
Given an arbitrary Boolean circuit, deciding whether its multilinear ANF has degree at most three is coNP-complete.
\end{proposition}
\begin{proof}
The complement has a polynomial certificate: four distinct coordinate directions and a Boolean point where the fourth partial difference is one. It is checked with 16 circuit evaluations. If the degree is at least four, differentiate a nonzero monomial in four of its variables; uniqueness of ANF coefficients implies that derivative is a nonzero Boolean function and has a witnessing point.

For hardness, map a circuit $C(x)$ to $G(x,z)=C(x)z_1z_2z_3z_4$, using four fresh input bits. If $C$ is unsatisfiable, $G$ is zero. Otherwise its nonzero ANF is multiplied by a disjoint monomial of degree four, so $\deg G\ge4$. This is a polynomial reduction from circuit unsatisfiability.
\end{proof}
A complete polynomial-time degree-acquisition stage for arbitrary circuits would therefore imply $\mathrm P=\mathrm{NP}$. The governing subexponential fixed-width counting claim does not assert that collapse. This is a boundary on the proposed generic program interface, not a counting lower bound for the structural family or a proof that every other program representation must fail.

\section{Decisive failures and the exact remaining scope}
\subsection{The unique-decoding radius is substantive}
Take $d=10$ and, on one field coordinate, the polar form of $\Tr(X^3)$. Its map is $L(Y)=Y^2+Y^{2^9}$. After precomposition by squaring this is $Y^4+Y$, with kernel $\F_4$; hence the polar rank is eight. Split its four symplectic pairs into two groups. A received alternating form at the midpoint has rank-four errors relative to both the zero compatible form and this nonzero compatible form. Both can be added to the same top cubic.

The bandwidth-one unique radius is three, so the two valid rank-four decompositions lie just beyond it. The program constructs this midpoint and rejects it as having no bounded-rank lift. This establishes nonuniqueness of the planted representation beyond the guarantee. It does not establish hardness of counting: either valid decomposition, if acquired, would still reconstruct the same scalar function.

\subsection{Trace aliases and Frobenius endpoints}
At $d=2$, a nonzero homogeneous ordinary cubic can disappear under scalar trace: $X^3$ takes values in $\F_2$, so $\Tr_{\F_4/\F_2}(X^3)=0$. At $d=1$, $X^3=X$ as functions. The special triple-factor acquisition consequently requires $d\ge3$.

The diagonal quadratic refinement also requires its linear correction. For example, over $F_8$, $\Tr(X^3)$ is one at the binary basis vector $1$, whereas a canonical squarefree binary quadratic is zero on every individual basis vector. Their polar forms agree but their functions do not until corrected.

At the excluded half-period boundary $2s=d$, positive and negative Frobenius indices coincide. If $c\in\F_{2^s}$, the map $L(X)=cX^{2^s}$ may be nonzero alternating, while $\Tr_F(cX^{1+2^s})$ is identically zero. A coefficient $a$ with $a+a^{2^s}=c$ gives the appropriate refinement, but the distinct-index proof and decoder dimension used here do not cover that case. The stated strict inequality prevents an invalid extension.

\subsection{A field is not always present}
For the cyclic cubic $\sum_i x_ix_{i+1}x_{i+2}$ on $n\ge7$ cyclic indices, the third tensor has zero radical and centroid exactly $\F_2 I$. The source archive gives a direct coefficient proof and exact checks for $7\le n\le24$. The present field compression then has $d=1$ and no useful bandwidth radius. A generic field must not be manufactured by merely naming the original binary vector space as a field; compatibility with the tensor is the essential acquired condition.

The alternative stabilizer construction below handles this cyclic foreground with suitable quadratic backgrounds. This does not erase the field compiler's boundary; it shows why different acquired representations remain useful.

\subsection{Counting is the demonstrated capability}
The theorem computes total scalar output counts. It does not automatically give equally cheap counts under arbitrarily many binary prefix constraints, since those constraints can increase the field-variable count or require many Fourier characters. Moreover, a nonconstant scalar Boolean cubic already has substantial support by the Reed--Muller minimum-distance bound. The exact counting result should not be advertised as a new generic one-bit inversion or rejection-sampling breakthrough.

The 96-bit execution and smaller complete counts show that the proposed counter is real and preserves observables. They do not establish practical superiority over equally informed field, tensor, stabilizer or algebraic solvers. No such performance claim is made.

\section{Two complementary acquired interfaces}\label{sec:other}
These two constructions were completed alongside the field compiler. Their acquisition proofs and finite programs are unconditional. Their stated subexponential local-phase consequences use the preserved Astra phase extension of the pinned manuscript \cite{eth}, including its manufactured replacement forms and exact reconstruction. The finite programs do not execute that asymptotic phase engine.

Write $F(N,m)$ for the earlier fixed-locality phase-counting exponent. The preserved sufficient regime is
\[
 \log(2+m/N)=o(\log N)\quad\Longrightarrow\quad F(N,m)=o(N),
\]
with polynomial space and charged exact reconstruction. This dependency is kept separate from Theorem~\ref{thm:main}.

\subsection{Weighted sparse stabilizers for a dense quadratic background}
Let $B$ be raw alternating, $P$ a sum of $m$ fixed-locality Boolean monomials, and $q_B(x)=\sum_{i<j}B_{ij}x_ix_j$. Manufacture the graph code
\[
 \mathcal C_B=\{(a,Ba):a\in\F_2^n\}.
\]
Search words of joint support at most $w$, in nondecreasing $X$ cost $\wt(a)$, and apply linear-matroid greedy with total $X$ budget $A_0$. For a support $U$, supported words are obtained from the kernel of $B[U^c,U]$. The weighted guarantee requires streaming all combinations of that kernel, not just an arbitrary basis.

Suppose the selected independent family has rank $r$, total $X$ cost $A$, total $Z$ cost $D$, and defect $\delta=n-r$. Then $A\le A_0$, $D\le rw$, and acquisition costs
\[
 \poly(n)\,2^{O(w\log(en/w))}\qquad(1\le w\le n/2).
\]
If an undisclosed independent family has rank $n-\delta_*$, total $X$ cost at most $A_0$, and total joint weight $W_*$, weighted greedy guarantees
\[
 \delta\le\delta_*+\left\lfloor\frac{W_*}{w+1}\right\rfloor.
\]
Deleting heavier words from that hidden family proves the bound; matroid greedy supplies the minimum cost for every attainable prefix rank.

The exact stabilizer identity is
\[
 K(a)|x\rangle=(-1)^{q_B(a)+(Ba)\cdot x}|x+a\rangle,
 \qquad K(a)K(b)=K(a+b).
\]
Complete the acquired vectors to a basis. The dense complement costs $2^\delta$ branches. For each acquired generator, a conditional $X$ update is enforced with one new state bit and one Fourier multiplier; a $Z$ action contributes a local phase and no state bit. This gives $2^\delta$ exact local phase sums, each with
\[
 N=r+2A,\qquad m'=m+3A+D+r+O(1),
\]
and common denominator $2^A$. All branch numerators are reconstructed before division. Their bit length is $O(n+A)$. The resulting conditional time is
\[
 \poly(L+n+A_0)\left[2^{O(w\log(en/w))}
 +2^{\delta+F(N,m')}\right].
\]
It is subexponential in the original $n$ if an undisclosed family has $A_0=O(n)$, $\delta_*=o(n)$, $W_*=n^{1+o(1)}$ and $m=n^{1+o(1)}$.

A covered family is $B=Q^{-1}+S$ with unknown invertible alternating $Q$ having $O(n)$ entries and maximum degree $n^{o(1)}$, and unknown symmetric sparse $S$ with $n^{1+o(1)}$ edges. Rows of $Q$ give the hidden generators; their total $X$ cost is $O(n)$ and total $Z$ cost is at most $n+2\Delta e(S)$. The inverse of the corrupted $B$ need not remain sparse. The result accommodates the cyclic cubic foreground, whose primitive tensor block has dimension $n$, so it removes the earlier small-primitive-block condition on this family.

The general boundary is decisive: a union bound shows that a random alternating $B$ has no nonzero graph-code word of joint support at most $n/8$ with probability tending to one exponentially fast. Small-support stabilizers cannot be assumed for every dense quadratic. Graph-code invariants, stabilizer identities and weighted-basis optimization are prior machinery; the scoped contribution is their acquired interface and variable accounting. Full proof, literature comparisons and execution of 278,796 pointwise trace identities are in \texttt{work\_gaussian\_new\_agent}.

\subsection{A polynomial annihilator manufactures a hidden low-rank code}
Let a raw symmetric matrix satisfy $B=D+S$, where $D=UKU^\top$ has rank $r$, $K$ is nonsingular symmetric, and each row of $S$ has weight at most $\Delta$. Suppose a monic binary polynomial $h$ of degree $q$ annihilates $S$. Define the unobserved orbit matrix
\[
 W=[U,SU,\ldots,S^{q-1}U].
\]
Assume $W$ has full rank $qr$ and its code has minimum distance greater than $2\Delta$. Neither $W$, $U$, $D$ nor $S$ is supplied.

Expanding $h(S+D)$ as noncommutative words gives
\[
 h(B)=WHW^\top,
\]
where in $r\times r$ blocks $H_{ij}=0$ for $i+j>q-1$ and $H_{ij}=K$ for $i+j=q-1$. The only terms attaining that outer degree are the one-$D$ words in the leading power. Reversing block columns makes $H$ triangular with nonsingular diagonal blocks $K$, so
\[
 \im h(B)=\im W,\qquad \rank h(B)=qr.
\]
This manufactures the needed code from the raw matrix. Enumerate it in polynomial space, verify its distance, and decode each column of $B$ within Hamming radius $\Delta$. The true $D$ column is the unique codeword in that radius. Verify symmetry, exact decomposition and $h(S)=0$. Acquisition costs $2^{qr}\poly(n+q)$, without a condition on rank times the number of sparse edges. For fixed $q$, a portfolio of all monic polynomials acquires $h$ too; code-rank budgets prevent an unfavorable trial from dominating the search.

With $M=e(S)$ and the preserved exact Gaussian refinement (including diagonal/unary corrections), the conditional original-input phase cost is
\[
 \poly(L+n+q)\left(2^{qr}+2^{r+F(n,m+M+O(n))}\right).
\]
For a hidden vertex permutation of the $\log_2n$-dimensional hypercube, $S^2=(\log_2n\bmod2)I$, so $q=2$ and $M=n\log_2n/2$. A random $n\times r$ matrix $U$, with even $r=o(n)$, has the required orbit rank and distance with high probability; the source proof gives an explicit union bound. Taking $r\asymp n/\log n$ gives $rM\asymp n^2/2$, outside the earlier sufficient condition $rM=o(n^2)$, while this acquisition still costs $2^{o(n)}$.

Transversality and distance are substantive restrictions. A commuting rank-two perturbation of a matching can make $B^2+I=0$ and erase the perturbation code entirely. Full orbit rank can also coexist with distance one. Both failures are executed. Ordinary annihilator identities are not invariant under arbitrary congruence, so this result is not silently extended through arbitrary binary gauges.

Ten fixtures recovered exact hidden matrices, including a 512-variable case with $r=6$, $q=2$, code dimension 12, distance 219 and 2,304 sparse edges. The tests checked 20,000 quadratic identities and 69,632 complete mixed-phase assignments. A separate coding argument and 12,094 executed distance checks show why choosing any sublinear number of noisy raw columns can miss the required code even with only matching errors. Full proofs, comparators and data are in \texttt{work\_rank\_next\_agent}.

\section{Prior machinery, comparators, and the defensible claim}
The ingredients are substantial prior algorithms. Tensor centroid descent already acquires hidden scalar fields \cite{magma}. Gabidulin and Loidreau decoding already correct rank errors within the stated radius \cite{silva,bombar}. Gaussian character expansions and Dwork point counting are established \cite{lauderwan}. When a field base and low-rank defect are already supplied, the Walsh-plus-point-counting combination is an equally informed comparator with the same $2^r$ leaf count. The question addressed here is how to acquire that interface from one raw scalar cubic and verify every term in its reconstruction.

Hidden-field perturbation recovery also has direct cryptographic ancestry. Internal HFE perturbations were already studied by Ding--Schmidt and Dubois--Granboulan--Stern \cite{ding,dubois}. Hat-plus uses a small output span of arbitrary quadratic perturbations \cite{hat}. The especially close HFE-LL' construction adds products of binary linear forms, and Smith-Tone--Valenzuela explicitly recover their noise support \cite{llprime}. A scalar quadratic of polar rank $r$ is, modulo affine terms, a sum of $r/2$ such products. Neither that perturbation shape nor the general idea of recovering it is new here.

The narrower candidate contribution is the deterministic, adversarial-error \emph{scalar exact-counting interface}: the untouched top tensor acquires the field, the trace pairing exposes a guaranteed rank-metric code, decoding recovers a bounded ordinary-degree base, and a real Fourier expansion preserves the original count at a fully charged cost. The bandwidth version and the exact gate front end extend the interface in two distinct directions. The inspected sources do not state this specific counting theorem, but the search is not a proof of historical priority. The official HFE-LL' abstract and indexed model material were available; the complete PDF retrieval was blocked, so a full algorithm-by-algorithm audit of that paper remains unresolved.

The other interfaces address different acquisition obstructions. Weighted stabilizers permit a large primitive nonlinear block; polynomial annihilators remove an earlier rank-times-sparsity condition on a concrete family. They do not make the field model generic. Likewise, finite experiments verify exact construction identities and real executable counters, not dominance over optimized competitors. No practical-speed claim or broad computational-complexity collapse is inferred.

\section{Executable files and reproducibility}
The source archive preserves the following ordinary programs and their actual output data. It requires Python 3 and NumPy for the implemented Dwork backend. No external solver, hidden field oracle, or exponential truth-table counter is used by the public field counter. Exhaustive oracles occur only in separately named validation routines.

The main modules are \path{work_field_root/field_kloosterman.py}, \path{field_descent.py}, \path{field_rank_decoder.py}, and \path{field_bandwidth.py}. The gate interface is \path{circuit_field_counter.py}. The command-line wrapper accepts a JSON object containing $n$ and a list of Boolean monomial masks. Its count method uses the complete triple-product backend; its acquisition method emits a verified wider representation without claiming to execute a general field counter.

From the extracted archive root:
\begin{verbatim}
python work_field_root/run_field_counter.py \
  work_field_root/example_21_bit_input.json \
  --output exact_count.json

python work_field_root/circuit_field_counter.py \
  work_field_root/example_36_bit_circuit.json \
  --output circuit_count.json

python work_field_root/run_field_counter.py \
  work_field_root/example_21_bit_input.json \
  --method acquire --bandwidth auto --output representation.json
\end{verbatim}
The first command returns $Z=128$, $N_0=1\,048\,640$; the gate command returns $Z=-176\,128$, $N_0=34\,359\,650\,304$. To reproduce the finite checks, run the \path{field_rank_checks.py}, \path{field_rank_failures.py}, \path{field_bandwidth_checks.py} and \path{circuit_boundary_checks.py} scripts. The 96-bit fixture generator and complete execution are preserved separately in \path{field_descent_96.py} and its results. Independent algebra and gate audits are in \path{work_program_next_agent}; the broader bandwidth audit is in \path{work_rank_next_agent}.

The archived proofs state exactly which results use the governing phase theorem. Earlier response-type and passive-network artifacts, including the preserved task-directed optimization source, are left unchanged. Their contents are not silently replaced by this field-family pivot.

{\small
\begin{thebibliography}{9}
\setlength{\itemsep}{1pt}
\setlength{\parskip}{0pt}
\bibitem{lauderwan}
A.~G.~B. Lauder and D. Wan.
\emph{Counting points on varieties over finite fields of small characteristic}.
Especially the splitting function, finite trace algorithm, Proposition 36 and Section 6.4; Theorem 37 computes the larger full-zeta object.
\url{https://arxiv.org/pdf/math/0612147}.

\bibitem{silva}
D. Silva and F.~R. Kschischang.
\emph{Fast Encoding and Decoding of Gabidulin Codes}, 2009.
\url{https://arxiv.org/pdf/0901.2483}.

\bibitem{bombar}
M. Bombar and A. Couvreur.
\emph{Right-hand side decoding of Gabidulin codes and applications}, 2021.
Preliminaries include the trace adjoint and the established left-annihilator formulation.
\url{https://arxiv.org/pdf/2112.07322}.

\bibitem{magma}
Magma handbook, \emph{Linear Invariants of Tensors},
\texttt{TensorOverCentroid} and its hidden-field worked example.
The documented ancestry includes P.~A. Brooksbank and J.~B. Wilson,
\emph{The module isomorphism problem reconsidered}, J. Algebra 421 (2015),
and U. First, J. Maglione and J.~B. Wilson's tensor operator theory.
\url{https://magma.maths.usyd.edu.au/magma/handbook/text/672}.

\bibitem{ding}
J. Ding and D. Schmidt.
\emph{Cryptanalysis of HFEv and Internal Perturbation of HFE}, PKC 2005.
\url{https://archive.ymsc.tsinghua.edu.cn/pacm_download/672/12592-dingjt6.pdf}.

\bibitem{dubois}
V. Dubois, L. Granboulan and J. Stern.
\emph{Cryptanalysis of HFE with Internal Perturbation}, PKC 2007.
\url{https://www.di.ens.fr/~stern/data/St118.pdf}.

\bibitem{hat}
J.-C. Faugere, G. Macario-Rat, J. Patarin and L. Perret.
\emph{A New Perturbation for Multivariate Public Key Schemes such as HFE and UOV}, 2022.
\url{https://vox-sign.com/files/hat.pdf}.

\bibitem{llprime}
D. Smith-Tone and C. Valenzuela.
\emph{Cryptanalysis of the best HFE-LL' Constructions}, 2025.
Official abstract and metadata inspected; the complete PDF was not retrieved.
\url{https://eprint.iacr.org/2025/1362}.

\bibitem{eth}
OpenAI mathematics release, user-pinned manuscript,
\emph{The Exponential Time Hypothesis is False}, 25 October 2026.
Commit \texttt{7ef5c0cc5ca4b0e072b7a050fca09ab24bc6fa3a}.
The relevant replacement and reconstruction machinery is in Section 4 and Proposition 3.4, with exact arithmetic in Appendix B. This reference is used only for the explicitly conditional local-phase consequences, not the field theorem.
\href{https://github.com/openai/math/blob/7ef5c0cc5ca4b0e072b7a050fca09ab24bc6fa3a/preprints/The-Exponential-Time-Hypothesis-is-False-October-25-2026/paper.pdf}{Pinned PDF at the specified commit}.
\end{thebibliography}
}
\end{document}
